% Note d'accompagnement des releves de notes - Candidature Bell 2027
% Identite typographique alignee sur le CV et la lettre : Carlito
% Palette : bleu #263068 et rouge #E6332A.
% Compilation : xelatex Note_accompagnement_releves_Bell_2027.tex
\documentclass[10pt,a4paper]{article}
\usepackage[a4paper,top=1.75cm,bottom=1.8cm,left=2.05cm,right=2.05cm]{geometry}
\usepackage{fontspec}
\setmainfont{Carlito}[BoldFont={Carlito Bold},ItalicFont={Carlito Italic},Ligatures=NoCommon]
\usepackage[french]{babel}
\usepackage{microtype}
\usepackage{xcolor}
\usepackage{setspace}
\usepackage{ragged2e}
\usepackage[hidelinks]{hyperref}
\definecolor{bleu}{HTML}{263068}
\definecolor{rouge}{HTML}{E6332A}
\definecolor{gris}{HTML}{9EA6B6}
\definecolor{anthracite}{HTML}{222936}
\color{anthracite}
\hypersetup{pdftitle={Note d'accompagnement - Releves de notes - Bell 2027},pdfauthor={Youssef EL KATTOUFI}}
\pagestyle{empty}
\setlength{\parindent}{0pt}
\setlength{\parskip}{11pt}
\setstretch{1.14}
\begin{document}
{\fontsize{15.4}{17.2}\selectfont\bfseries\color{bleu}Youssef EL KATTOUFI}\par\vspace{2pt}
{\fontsize{10.1}{11.4}\selectfont\bfseries\color{rouge}Élève-ingénieur -- Électronique et informatique}\par\vspace{1pt}
{\fontsize{13}{15}\selectfont\bfseries\color{bleu}Note d'accompagnement -- Relevés de notes}\par
\vspace{11pt}
\fontsize{11}{14.6}\selectfont
Madame, Monsieur,

Ne sachant pas précisément quels relevés de notes étaient attendus pour ma candidature, j'ai préféré joindre \emph{l'ensemble de mes relevés depuis le début de mes études supérieures}, afin de vous donner une vue complète de mon parcours.

Si vous souhaitez consulter les relevés correspondant à ma formation actuelle à \textbf{Polytech Sorbonne}, ils se trouvent \textbf{à la fin du dossier}. Il suffit de faire défiler le document jusqu'aux pages portant le logo et la mention « Polytech Sorbonne ».

Je vous remercie pour l'attention portée à ma candidature.
\end{document}
